% BÀI TOÁN HÌNH HỌC - BAI_HINH_52 (TRÍCH ĐỀ 52)
\begin{ex}
Cho đường tròn $(O; R)$ và dây cung $AB$ khác đường kính. Gọi $H$ là trung điểm của $AB$, đường thẳng $HO$ cắt đường tròn $(O)$ tại hai điểm $C$ và $D$ ($C$ thuộc cung nhỏ $AB$). Điểm $M$ thuộc đoạn $AH$ ($M$ khác $A$ và $H$), tia $CM$ cắt đường tròn $(O)$ tại điểm thứ hai là $N$.
	\begin{enumerate}
		\item[a)] Chứng minh 4 điểm $D, H, M, N$ nằm trên một đường tròn.
		\item[b)] Gọi $I$ là giao điểm của hai đường thẳng $DN$ và $AB$, $K$ là giao điểm của hai đường thẳng $CI$ và $DM$. Chứng minh $NM$ là đường phân giác trong của góc $\widehat{KNH}$ và $M$ là tâm đường tròn nội tiếp tam giác $HKN$.
	\end{enumerate}
\end{ex}
